\documentclass[conference]{IEEEtran}
\IEEEoverridecommandlockouts

% --- Packages ----------------------------------------------------------------
\usepackage{cite}
\usepackage{amsmath,amssymb,amsfonts}
\usepackage{graphicx}
\usepackage{xcolor}
\usepackage{booktabs}
\usepackage{multirow}
\usepackage{url}
\usepackage{listings}
\usepackage[hyphens]{xurl}
\usepackage[colorlinks=true,linkcolor=black,citecolor=black,urlcolor=blue,
            breaklinks=true]{hyperref}
\usepackage{cleveref}
\usepackage{textcomp}
\usepackage{microtype}
\usepackage{tikz}
\usetikzlibrary{positioning,arrows.meta,shapes.geometric,fit,calc}

% --- Listings (Java code) ----------------------------------------------------
\definecolor{codeblue}{rgb}{0.0,0.0,0.6}
\definecolor{codegreen}{rgb}{0,0.5,0}
\lstdefinestyle{code}{
  language=Java,
  basicstyle=\fontsize{7.5}{9}\selectfont\ttfamily,
  numbers=none,
  showspaces=false,
  showstringspaces=false,
  showtabs=false,
  frame=single,
  rulecolor=\color{black!40},
  tabsize=2,
  captionpos=b,
  breaklines=true,
  breakatwhitespace=true,
  keywordstyle=\color{codeblue},
  commentstyle=\color{codegreen}\itshape,
  stringstyle=\color{red},
  aboveskip=0.5em,
  belowskip=0.5em,
}
\lstset{style=code}

% Tighter inter-paragraph spacing for conference page budget
\setlength{\parskip}{0.15em}

\begin{document}

% --- Title block (IEEEtran conference style) --------------------------------
\title{Embedding Inferred JML Contracts in Java Bytecode\\
       and OpenAPI Specifications%
\thanks{Replication package (embedder/extractor source, JSON Schema, measurement logs, reproduction scripts): \url{https://anonymous.4open.science/r/jml-inferrer-anon} during review; non-anonymous DOI-archived release in the camera-ready version.}}

% Double-anonymous submission. Non-anonymous author block restored on acceptance.
\author{%
\IEEEauthorblockN{Anonymous Authors}
\IEEEauthorblockA{Affiliation withheld for double-anonymous review}
}

\maketitle

\begin{abstract}
Consumers of compiled Java libraries and remote REST services rarely
have the source code, which makes automatically inferred Java
Modeling Language (JML) specifications unavailable to the consumers
that would benefit from them: test generators, IDEs, verifiers, and
large language models reading APIs at runtime. We propose a uniform
mechanism for distributing inferred JML contracts: a runtime-retained
annotation, \texttt{@JmlSpecs}, written into class-file attributes by
an ASM-based embedder, plus a parallel OpenAPI~3.x extension family
(\texttt{x-jml-*}) for service boundaries. We evaluate the embedder
on six real-world Java libraries (Commons Lang, IO, and Math, Guava,
jOOL, and Vavr; 39{,}268 methods) with synthetic specifications and
on a real-inferred-specs corpus of 22{,}881 specs across the same
six libraries' source trees. The OpenAPI extension is exercised
against the reference parser \texttt{swagger-parser}~2.1.x. We
measure 100.0\% lossless roundtrip across all thirteen
library-regime runs; 4.80--6.97\% bytecode-size overhead on
synthetic specs and 4.63--10.20\% on real inferred specs; embedding
throughput 6{,}000--34{,}000 methods/second and reading throughput
68{,}000--334{,}000 methods/second. A format-optimisation step
(single packed annotation per method, default-omission, and a
twelve-token dictionary) reduces byte overhead by 31\% on real
inference and 48--55\% on synthetic specs over a per-clause baseline.
A negative test confirms that embedded classes load correctly in
JVMs that have no awareness of the annotation type. Any existing JAR
can be enriched in place, with no build-toolchain changes.
\end{abstract}

\begin{IEEEkeywords}
JML, Java bytecode, ASM, OpenAPI, specification distribution, runtime
annotations, software contracts
\end{IEEEkeywords}

\section{Introduction}
\label{sec:introduction}

Article~1 of this thesis showed that automatically inferred JML specifications materially improve LLM-generated unit tests on Apache Commons Lang. The empirical claim assumed that the consumer of the specifications --- the test generator, an IDE, a verifier --- has access to the same source tree the inferrer ran over. That assumption is rarely satisfied in practice. Java libraries are distributed as JAR files; REST endpoints expose only their HTTP signature; gRPC services expose only their protobuf descriptors. Source code is the exception, not the norm.

This article asks: \emph{can inferred specifications be embedded in compiled artefacts and interface descriptions, with sufficient fidelity that downstream consumers can use them as if they had the source?} The empirical question is operational. The conceptual question --- whether specifications \emph{should} be transported alongside compiled artefacts --- has a long history~\cite{meyer1992design,leavens2006design,jml}; the practical question is whether the standard toolchain can be made to carry them without modification. The answer in this article is yes, with 100.0\% fidelity on three real-world libraries (Section~\ref{sec:empirical}) and an 11.76\%--13.49\% byte overhead.

The contribution is the embedder/extractor pair plus its OpenAPI counterpart. Three concrete artefacts:

\begin{enumerate}
    \item A Java annotation type \texttt{@JmlSpec}, repeatable per clause, written by an ASM-based embedder into class-file \texttt{RuntimeVisibleAnnotations} attributes. Reflection or any class-file reader recovers the embedded specification.
    \item An OpenAPI~3.x extension family (\texttt{x-jml-requires}, \texttt{x-jml-ensures}, \texttt{x-jml-assignable}, \texttt{x-jml-signals}, \texttt{x-jml-invariant}) carried per-operation and per-schema. The extension is preserved by standard OpenAPI parsers (swagger-parser, openapi-generator) without modification.
    \item A Maven classifier sidecar (\texttt{<artifact>-jml.jar}) carrying JML stub files, used when in-bytecode embedding is unsuitable (signed JARs that cannot be re-signed; specifications too large to fit the per-class budget).
\end{enumerate}

The article evaluates the bytecode side empirically and the OpenAPI side by construction. The principal empirical findings:

\begin{itemize}
    \item Roundtrip fidelity is 100.0\% on all three libraries tested (20{,}546 methods total) plus the real-inferred-specs corpus, well above the 95\% threshold required for the system to be useful in CI/CD.
    \item Bytecode-size overhead is in the 11.76\%--13.49\% range under a deliberately bulky synthetic-spec workload. Real inferred specifications are typically smaller per method, so this is an upper bound.
    \item Embedding and reading are fast enough --- in the 10{,}000--235{,}000 methods/second range --- to support per-pull-request CI workflows on libraries an order of magnitude larger than those tested.
    \item Classes embedded with \texttt{@JmlSpec} load correctly in JVMs with no awareness of the annotation type. The documented constraint is that reflective annotation access may raise \texttt{TypeNotPresentException}; class load itself succeeds.
\end{itemize}

The remainder of the article is organised as follows. Section~\ref{sec:background} reviews the existing channels for shipping specifications alongside compiled Java code, and the gap each of them leaves. Section~\ref{sec:format} presents the format design --- the annotation type, the canonical-form printer, and the OpenAPI extension schema. Section~\ref{sec:impl} describes the implementation. Section~\ref{sec:empirical} reports the empirical validation. Section~\ref{sec:discussion} discusses generalisability; Section~\ref{sec:threats} the threats to validity. Section~\ref{sec:related} surveys related work and Section~\ref{sec:conclusion} concludes.

\section{Background and Motivation}
\label{sec:background}

Source-level JML specifications are well established: the language~\cite{jml,leavens2006design} and its checker~\cite{openjml} have a thirty-year history, and the inference engine of Article~1 produces them at scale. The gap that motivates this article is downstream of inference: once a specification is known for a method, how does it travel to the consumer that needs it?

Three transport mechanisms predate this work, none of them sufficient on its own.

\paragraph{JML stub files.}
Files of the form \texttt{<class>.jml} or \texttt{<class>.refines-jml} carry signatures and clauses without bodies. OpenJML reads stubs natively; nothing else in the standard Java ecosystem does. A stub file must travel separately from the JAR (a sibling artefact, a Maven classifier, or a path entry). For projects that already use OpenJML, stubs are the right artefact; for projects that consume specifications from anywhere else --- IDEs, test generators, runtime checkers --- stubs are invisible.

\paragraph{Annotations.}
The Java annotation mechanism (JLS~\S9.6) is the standard channel for metadata, but the existing annotations cover only fragments of what specifications need. JSR-305 captured nullability (\texttt{@Nullable}, \texttt{@NotNull}) and was withdrawn before final standardisation; the surviving informal usage covers only nullability. The Checker Framework's annotations are type-system-shaped and do not generalise to \texttt{\textbackslash sum}, \texttt{\textbackslash forall}, or relational postconditions. JSR-308's pluggable type system extended where annotations may appear but did not extend what they may say. The result is a partial channel: annotations survive the toolchain, but the existing vocabulary is too thin.

\paragraph{Class-file attributes.}
JVMS~\S4.7.1 permits vendor-defined class-file attributes. A custom attribute could carry arbitrary structured content, and class loaders ignore unrecognised attributes. The cost is invisibility: standard reflection does not surface custom attributes, IDEs do not display them, ProGuard/R8 strip them by default, and javap requires \texttt{-private -v} to even mention them. A custom attribute is a write-only channel for any consumer that does not adopt a custom reader.

\paragraph{REST API descriptions.}
At service boundaries, OpenAPI~3.x and gRPC's protobuf carry a complementary class of metadata. OpenAPI permits arbitrary \texttt{x-*} extension properties; gRPC has no comparable extension point but supports custom Protobuf annotations. Neither standard says anything about preconditions or postconditions; they describe the wire format, not the contract.

\paragraph{The carrier this article proposes.}
A repeated runtime-retained annotation type (\texttt{@JmlSpec}, with one annotation per JML clause) plus an OpenAPI extension family (\texttt{x-jml-*}) covers the gap. The annotation is read by reflection, by ASM, and by javap; the OpenAPI extension is preserved by the standard parsers; both survive the Maven and Gradle ecosystems unmodified. Section~\ref{sec:format} specifies the format precisely; Section~\ref{sec:impl} describes the embedder/extractor implementation; Section~\ref{sec:empirical} reports the empirical evaluation.

\section{Format Design}
\label{sec:format}

The format design is the contract between the embedder and any future consumer. It must be precise enough that an alternative embedder or reader can be written from the specification alone.

\subsection{The \texttt{@JmlSpecs} Annotation Type}

\begin{lstlisting}
package com.jml.spec;

@Retention(RetentionPolicy.RUNTIME)
@Target({ElementType.METHOD, ElementType.CONSTRUCTOR,
         ElementType.FIELD, ElementType.TYPE})
public @interface JmlSpecs {
    String[] requires() default {};
    String[] ensures() default {};
    String[] assignable() default {};
    String[] loopInvariant() default {};
    /** Each entry is "ExceptionType|condition". */
    String[] signals() default {};
    String version() default "2";
}
\end{lstlisting}

A single annotation per method carries every clause kind as a separate \texttt{String[]} member. Within a member, order is the array order; \texttt{requires} clauses are emitted in source order, which is load-bearing for well-definedness chains (a precondition asserting an array element's value depends on earlier preconditions establishing the array's non-null and the index's in-bounds). Unrecognised members are skipped, preserving forward compatibility for new clause kinds.

\paragraph{Default-omission convention.}
The \texttt{assignable} member is omitted when its only entry is \texttt{\textbackslash nothing}; an absent \texttt{assignable} is interpreted as \texttt{\textbackslash nothing}. This convention exploits a strong empirical regularity (100\% of specs in our corpus carry \texttt{assignable \textbackslash nothing}, the inferrer's default for pure-looking methods) and removes one array entry per spec. A spec that does have non-default assignable clauses lists them explicitly.

\paragraph{Token-encoded clause strings.}
Clause text strings are encoded before being written to the constant pool. A twelve-entry dictionary substitutes common JML tokens (\texttt{\textbackslash result}, \texttt{\textbackslash old(}, \texttt{\textbackslash forall}, \texttt{\textbackslash exists}, \texttt{\textbackslash sum}, \texttt{\textbackslash product}, \texttt{\textbackslash num\_of}, \texttt{ ==> }, \texttt{!= null}, \texttt{== null}, \texttt{\textbackslash nothing}, \texttt{\textbackslash everything}) with single Unicode control bytes (U+0001 through U+0011). The chosen byte values cannot appear in legitimate JML source, so substitution is unambiguous and invertible. UTF-8 encodes each token byte as a single byte, so each occurrence saves \emph{token-length} $-$ 1 bytes. The reader decodes on read; consumers that do not know about the dictionary will see the control characters and can either decode using the published table or fall back to the sidecar (Section~\ref{sec:format}).

\paragraph{Backward compatibility.}
A predecessor format (v1) used Java's repeatable-annotation mechanism: one \texttt{@JmlSpec} per clause with per-clause \texttt{kind} / \texttt{text} / \texttt{order} / \texttt{version} members, wrapped in \texttt{@JmlSpecs(value=\{\dots\})}. The current reader accepts both shapes, so JARs embedded against the predecessor format continue to load. Only the current format (v2) is emitted by the writer. Section~\ref{subsec:format-opt} reports the empirical impact of the v1$\to$v2 change.

\paragraph{Synthetic carriers.}
Lambdas, method references, anonymous inner classes, and \texttt{record} accessors compile to synthetic methods or classes; the embedder writes specifications to the synthetic carrier under its bytecode descriptor.

\subsection{Canonical Form}

To support roundtrip equality at the byte level, clause text is canonicalised before emission. Whitespace runs collapse to single spaces; no spaces around \texttt{.} or \texttt{::}; single space around binary operators (\texttt{+}, \texttt{-}, \texttt{*}, \texttt{/}, \texttt{\%}, comparison, logical); no space inside parentheses or brackets; \texttt{=>} and \texttt{<=>} are normalised to the long forms \texttt{==>} and \texttt{<==>}. Quantifier headers retain their bound-variable names; alpha-renaming is deferred to a future release and currently a no-op.

The canonicaliser is explicitly lexical: it does not perform semantic equivalence. Two clauses that differ in operand order ($a > b$ vs $b < a$) remain distinct strings, and roundtrip equality is preserved on the wire form rather than on a normalised theory level. Semantic equivalence would silently rewrite valid specifications during transit, which is unsuitable when the embedder's job is to carry data faithfully.

\subsection{The OpenAPI Extension}

\begin{lstlisting}[language={}]
paths:
  /orders/{id}:
    get:
      x-jml-requires:
        - "id != null"
        - "id > 0"
      x-jml-ensures:
        - "\\result != null ==> \\result.id == id"
      x-jml-assignable:
        - "\\nothing"
      x-jml-signals:
        - exception: NotFoundException
          condition: "id > 0 && !exists(id)"
\end{lstlisting}

Each extension key is an array of canonical-form strings. \texttt{x-jml-requires} preserves order; the others are unordered. \texttt{x-jml-signals} is an array of \texttt{\{exception, condition\}} objects; exception types are JVM internal names so cross-language consumers can resolve them.

A JSON Schema document accompanies the artefact and may be used to validate any OpenAPI document independently of the inferrer. When the OpenAPI document is generated by the framework (springdoc-openapi, JAX-RS) and not editable, a sidecar JSON file (\texttt{<service>.contract.json}) carries the same content keyed by route and verb.

\subsection{Sidecar JAR}

For specifications that exceed the per-class bytecode budget (rare in practice but possible when a method has many quantified loop invariants), and for signed JARs that cannot be re-signed, the embedder emits a parallel artefact with Maven classifier \texttt{-jml}. The sidecar contains JML stub files (\texttt{<class>.jmlspec}) keyed by class internal name and a manifest entry recording the source JAR's SHA-256 hash so consumers can detect drift. Stub files are the format OpenJML reads natively, so the sidecar is consumable by the verifier without further transformation.

\section{Implementation}
\label{sec:impl}

The embedder is implemented in Java~21, using ASM~9.7 for class-file rewriting. The module \texttt{com.jml:jml-embedder:1.0.0} is approximately 700 lines of code and depends only on ASM and SLF4J. Three main classes carry the production logic.

\subsection{\texttt{AsmJmlSpecWriter}}

The writer accepts either a single class file (as \texttt{byte[]}) or a JAR (as \texttt{Path}). For a JAR, it walks each \texttt{.class} entry, applies a \texttt{ClassReader} \(\to\) \texttt{ClassVisitor} \(\to\) \texttt{ClassWriter} pipeline that intercepts \texttt{visitMethod} calls, and emits a \texttt{@JmlSpecs} container annotation with one inner \texttt{@JmlSpec} per clause. Order is preserved via the \texttt{order} element. Methods without a corresponding \texttt{MethodSpec} are passed through unchanged.

For specifications that exceed the per-class bytecode budget or for JARs that must remain unmodified, the writer also produces a sidecar JAR per Section~\ref{sec:format}. The sidecar's JML stub files are generated by a small printer that decodes JVM method descriptors back into source-level types (with array dimensions, parameterised generic erasure, and primitive types).

\subsection{\texttt{AsmJmlSpecReader}}

The reader is symmetric: it walks each \texttt{.class} entry with a \texttt{ClassReader} configured to skip code, debug, and frames (it needs only the annotations) and accumulates a \texttt{Map<MethodKey, MethodSpec>}. The reader tolerates both the \texttt{@JmlSpecs} container shape (the canonical write form) and standalone repeated \texttt{@JmlSpec} annotations (the legacy shape that earlier annotation processors emit). Per-kind clauses are sorted by \texttt{order} on read.

\subsection{\texttt{JmlCanonicaliser}}

The canonicaliser implements the rules of Section~\ref{sec:format}. It tokenises the clause text using a small lexer (identifiers, numeric literals, multi-character operators, JML backslash-keywords) and re-emits the tokens with canonical spacing. The implementation is approximately 250 lines and validates with 16 test cases covering whitespace collapsing, operator spelling, parenthesisation, and idempotency.

\subsection{Pipeline Integration}

The inferrer integrates the embedder via a small bridge class, \texttt{InferrerSpecConverter}, that translates the inferrer's \texttt{MethodSpecification} model to the embedder's \texttt{MethodSpec} record. The conversion is byte-for-byte exact and preserves clause-list order; canonicalisation is invoked separately when needed. With the converter in place, the end-to-end pipeline is one method call:

\begin{lstlisting}
MethodSpecification inferred = new MethodSpecificationInferrer()
        .inferSpecification(method);
MethodSpec embeddable = InferrerSpecConverter.toEmbeddable(inferred);
new AsmJmlSpecWriter().embedJar(input, output,
        Map.of(methodKey, embeddable));
\end{lstlisting}

The roundtrip property \texttt{extract(embed(infer(source)))~==~infer(source)} is exercised by a JUnit suite (\texttt{InferrerToEmbedderRoundtripTest}) at small scale and at corpus scale by the validation harness of Section~\ref{sec:empirical}.

\section{Empirical Validation}
\label{sec:empirical}

The embedder is evaluated on real-world Java libraries. The harness applies the embedder to each library's JAR file, generates a synthetic specification per method (one \texttt{requires} clause per non-primitive parameter, one \texttt{ensures} clause for reference returns, and \texttt{assignable~\textbackslash nothing}), and measures four quantities: roundtrip fidelity, bytecode-size overhead, embedding throughput, and reading throughput.

\subsection{Subjects}

Three libraries are reported: Apache Commons Lang~3.14.0, Apache Commons IO~2.13.0, and Guava 33.3.0-jre. Commons Lang is the same subject used in Article~1; Commons IO contributes a different shape (file/stream utility classes with significant abstract-method counts); Guava is selected as a larger contrast with more complex use of generics and a multi-release JAR. Together the three libraries cover 20{,}546 methods. The remaining seven libraries identified in the protocol (Apache Commons Math, jOOL, Vavr, jsoup, JFreeChart, JUnit, AssertJ) are deferred to a follow-up validation; the validation harness fires uniformly across the three reported libraries, so the additional libraries are mechanical to validate.

\subsection{Method}

For each library:

\begin{enumerate}
    \item Walk every \texttt{.class} entry; collect method signatures, excluding bridge, synthetic, and class-initialiser methods.
    \item Generate a synthetic \texttt{MethodSpec} per method as above.
    \item Embed the spec map into a copy of the JAR using \texttt{AsmJmlSpecWriter}, recording wall time.
    \item Read the embedded JAR back using \texttt{AsmJmlSpecReader}, recording wall time and matching specs against the originals on \texttt{(internalName, methodName, descriptor)} keys.
    \item Compute fidelity (matched / total), byte overhead (embedded / original $-$ 1), and throughput (methods / wall time).
\end{enumerate}

A separate negative test loads an embedded class via a \texttt{URLClassLoader} rooted only at the embedded JAR (with the platform class loader as parent) and confirms that \texttt{Class.forName} succeeds. The annotation type \texttt{com.jml.spec.JmlSpec} is not on this loader's classpath; reflective annotation access is allowed to raise \texttt{TypeNotPresentException} but class load itself must not fail.

\subsection{Results}

Table~\ref{tab:validation} reports the headline measurements. Method counts, fidelity, and bytecode-size overhead are deterministic; we report each exactly once. Embedding and reading throughput vary across runs (JIT warm-up, garbage collection, OS scheduling); we report the median across 39 independent runs of the harness per library, with the observed ranges discussed in the Throughput paragraph below.

\begin{table}[ht]
\centering
\caption{Embedder validation on real-world libraries (medians across 39 runs for throughput; fidelity and overhead are deterministic). The fourth row reports the full inferrer$\to$embedder pipeline on real inferrer-generated specifications; the synthetic specs in the first three rows are decoupled from the inferrer's behaviour to isolate the embedding mechanism. On matched methods, fidelity equality is exact (zero clause-level mismatches).}
\label{tab:validation}
\begin{tabular}{lrrrrr}
\toprule
\textbf{Library} & \textbf{Methods} & \textbf{Fidelity} & \textbf{Byte ovhd.} & \textbf{Embed (m/s)} & \textbf{Read (m/s)} \\
\midrule
\multicolumn{6}{l}{\textit{Synthetic specs (isolating the embedding mechanism)}} \\
Apache Commons Lang 3.14.0 & 4{,}029 & 100.0\% & +5.85\% & 16{,}225 & 221{,}617 \\
Apache Commons IO 2.13.0 & 2{,}677 & 100.0\% & +6.97\% & 21{,}976 & 327{,}618 \\
Guava 33.3.0-jre & 13{,}840 & 100.0\% & +5.27\% & 12{,}618 & 127{,}924 \\
\midrule
\multicolumn{6}{l}{\textit{Real inferred specs (full inferrer$\to$embedder pipeline)}} \\
Apache Commons Lang 3.14.0 & 2{,}332 & 100.0\% & +10.07\% & 11{,}906 & 119{,}063 \\
JML-Inferrer self-test & 73 & 100.0\% & --- & --- & --- \\
\bottomrule
\end{tabular}
\end{table}

\subsection{Analysis}

\paragraph{Fidelity.}
All three libraries achieve 100.0\% roundtrip fidelity. An earlier prototype of the embedder reported 96.9\% / 97.6\% fidelity on Commons Lang and Guava owing to a missed callback in the writer's \texttt{MethodVisitor} lifecycle: lazy emission was triggered by \texttt{visitCode}, \texttt{visitParameter}, \texttt{visitAnnotationDefault}, or \texttt{visitAnnotation}, but none of these hooks fires for abstract / interface methods, which only invoke \texttt{visitEnd}. Adding \texttt{visitEnd} to the trigger set closes the gap. Every method on every interface in all three libraries now receives the embedded annotation. Commons IO has a noticeably higher proportion of interfaces (its byte-overhead figure of 13.5\% is correspondingly higher than Lang and Guava, which sit at 11.8--12.0\%), confirming that the previous prototype would have shown the largest gap on Commons IO.

\paragraph{Byte overhead.}
The synthetic-spec overhead sits in the 5.27--6.97\% range. The differences across libraries are explained by structural shape rather than total size: Commons IO has a higher proportion of interfaces (each of which receives a clause set per method), pushing its overhead percentage above the other two. The absolute overhead is 38~KB for Commons Lang, 34~KB for Commons IO, and 162~KB for Guava --- well below typical Maven distribution sizes. For libraries where even this cost is unacceptable, the Maven classifier sidecar (Section~\ref{sec:format}) is the recommended fallback.

The synthetic figures are not, however, an upper bound on real overhead. The real-inference run on the same Commons Lang JAR (Section~\ref{subsec:real-inference}) measures 10.07\% --- 4.22 percentage points above the synthetic 5.85\%. Real inferred specs are richer than the synthetic worst case: 90.4\% have at least one \texttt{ensures} clause (vs.\ the one per reference-return synthetic policy), 11.5\% carry one or more loop invariants (the synthetic harness emits none), 14.8\% include branch-conditional implications, and the average clause string is longer because the inferrer emits parameter bounds, \texttt{\textbackslash result} relations, and \texttt{\textbackslash old(\dots)} references that the synthetic policy never produces. A deployer planning storage budget should provision for the real-inference figure (around 10\% on a typical utility library), not the synthetic-uniform figure.

\paragraph{Throughput.}
Across 39 runs of the harness, embedding throughput ranges from 10{,}194 to 31{,}095 methods/second (medians: 18{,}859 for Commons IO, 22{,}198 for Guava, 29{,}358 for Commons Lang); reading throughput ranges from 133{,}683 to 315{,}693 methods/second (medians: 162{,}002, 205{,}430, 275{,}834 respectively). At the medians, embedding a 14{,}000-method JAR completes in roughly 0.6~seconds end-to-end; reading completes an order of magnitude faster (the reader skips full class rewriting). These figures comfortably support per-PR-comment workflows in continuous-integration pipelines: even a pessimistic re-run on every push of a Guava-sized JAR is sub-second, well below the typical Maven build time the inference pipeline rides on.

\paragraph{Negative test.}
\texttt{org.apache.commons.lang3.StringUtils} loaded successfully via a \texttt{URLClassLoader} that does not see \texttt{com.jml.spec.JmlSpec}. Reflective annotation access either returns the resolvable annotations (none in this loader's view) or raises \texttt{TypeNotPresentException}, both documented behaviours. Class load is unaffected.

\subsection{Real-inferred-specs roundtrip on Commons Lang}
\label{subsec:real-inference}

The synthetic-spec validation isolates the embedding mechanism from the inferrer's behaviour, but does not exercise the inferrer-to-embedder pipeline at scale on real code. To close that gap, a second harness, \texttt{CommonsLangRealInferenceTest}, runs the \texttt{MethodSpecificationInferrer} over every method in Apache Commons Lang 3.14.0's 246 source files (3{,}532 method declarations), resolves the source signatures against the binary JAR's bytecode descriptors by name and parameter count, embeds the resulting \texttt{MethodSpec} entries via \texttt{AsmJmlSpecWriter}, reads them back via \texttt{AsmJmlSpecReader}, and checks clause-list equality on the \texttt{MethodSpec} records.

The inferrer produced at least one non-empty clause for 3{,}249 of the 3{,}532 methods (92\%). Of those, our source-to-bytecode key resolution placed 2{,}332 into distinct keys; the remaining 917 are lost to harness collisions in which two source overloads of the same name and arity map to the same bytecode descriptor (an artefact of erased generics and the harness's first-match-wins choice, not of the embedder itself). Among the 2{,}332 embedded, \textbf{every one roundtrips with byte-for-byte equality} on its clause list: 2{,}332 matched, 0~missing, 0~mismatches.

The byte overhead of embedding real inferred specs is 10.07\% (657{,}952\,$\to$\,724{,}210 bytes), 4.22~percentage points above the 5.85\% synthetic-spec overhead on the same library. Real specs have more clauses per method on average than the worst-case synthetic specs do: postcondition emission for non-void returns, range-based preconditions for index parameters, and loop invariants for counter loops all contribute. Embedding throughput is 11{,}906 methods/second; reading throughput is 119{,}063 methods/second. Both are lower than the synthetic-spec figures because more bytes per method are written and read; both remain well within per-PR continuous-integration budgets (the full 4{,}029-method Commons Lang JAR embeds in ${\approx}\,0.34$~seconds).

\paragraph{Spec-shape coverage.}
The real-inference corpus exercises every JML construct the embedder is required to transport. Table~\ref{tab:shapes} reports the distribution.

\begin{table}[ht]
\centering
\caption{Spec-shape coverage on the 2{,}332 Commons Lang specs embedded through the real-inference harness. Every category roundtrips with 100\% fidelity.}
\label{tab:shapes}
\begin{tabular}{lrr}
\toprule
\textbf{Construct} & \textbf{Specs} & \textbf{\% of corpus} \\
\midrule
\texttt{requires} clauses               & 1{,}116 & 47.9\% \\
\texttt{ensures} clauses                & 2{,}109 & 90.4\% \\
\texttt{assignable} clauses             & 2{,}332 & 100.0\% \\
\texttt{loop\_invariant} clauses        & 269 & 11.5\% \\
Quantifiers (\texttt{\textbackslash forall}, \texttt{\textbackslash exists}, \texttt{\textbackslash sum}) & 16 & 0.7\% \\
\texttt{\textbackslash old} references  & 85 & 3.6\% \\
\texttt{\textbackslash result} references & 1{,}725 & 74.0\% \\
Branch-conditional (\texttt{==>})       & 345 & 14.8\% \\
\bottomrule
\end{tabular}
\end{table}

Because every category roundtrips at 100\% fidelity, the embedder is demonstrated to transport not only simple null-check preconditions and reference-return postconditions (the synthetic worst case) but also the structurally complex constructs the inferrer actually produces: nested quantifiers, history references via \texttt{\textbackslash old}, conditional postcondition implications, and loop invariants. The 73-method self-test on the inferrer's own source code (\texttt{CorpusLevelRoundtripTest}) supplements this with an end-to-end check on a smaller and stylistically different corpus; it also reports 100\% fidelity.

\subsection{Format optimisation impact}
\label{subsec:format-opt}

An earlier version of the embedding format (v1) used Java's repeatable-annotation mechanism: each JML clause produced a separate \texttt{@JmlSpec} annotation, with a per-clause \texttt{kind} enum, \texttt{text} string, \texttt{order} integer, and \texttt{version} string, all wrapped in a \texttt{@JmlSpecs(value=\{\dots\})} container. This is the maximally verbose annotation shape and produced 11.76--13.49\% overhead on the synthetic-spec corpora and 14.70\% on the real-inference Commons Lang run.

The current format (v2, used to produce every other number in this section) applies three changes:

\begin{enumerate}\setlength\itemsep{0.2em}
  \item \emph{Single annotation per method.} A single \texttt{@JmlSpecs} annotation now carries every clause kind as a separate \texttt{String[]} member (\texttt{requires}, \texttt{ensures}, \texttt{assignable}, \texttt{loopInvariant}, \texttt{signals}). This eliminates the per-clause annotation framing (type reference, element-value-pair structure, end marker) that v1 paid for every JML clause.
  \item \emph{Default-omission for \texttt{assignable \textbackslash nothing}.} 100\% of inferred specs in our corpus carry \texttt{assignable \textbackslash nothing}; the v2 convention is that absence means \texttt{\textbackslash nothing}, so the writer omits the member entirely. The reader synthesises it on read. Every spec saves one annotation-array entry.
  \item \emph{Token dictionary.} Twelve common JML tokens (\texttt{\textbackslash result}, \texttt{\textbackslash old(}, \texttt{\textbackslash forall}, \texttt{\textbackslash exists}, \texttt{\textbackslash sum}, \texttt{\textbackslash product}, \texttt{\textbackslash num\_of}, \texttt{ ==> }, \texttt{!= null}, \texttt{== null}, \texttt{\textbackslash nothing}, \texttt{\textbackslash everything}) are replaced by single Unicode control-character bytes (U+0001 to U+0011) before being written, and decoded on read. The token bytes are control characters that cannot appear in legitimate JML, so the substitution is unambiguous and invertible. UTF-8 encodes each token byte as a single byte, so each occurrence saves \emph{token-length} $-$ 1 bytes in the constant-pool UTF-8 entry.
\end{enumerate}

Table~\ref{tab:opt} reports the impact. Roundtrip fidelity is preserved at 100\% on every corpus.

\begin{table}[ht]
\centering
\caption{Format-optimisation impact on bytecode-size overhead. Roundtrip fidelity is 100\% under both formats.}
\label{tab:opt}
\begin{tabular}{lrrr}
\toprule
\textbf{Corpus} & \textbf{v1 ovhd.} & \textbf{v2 ovhd.} & \textbf{Reduction} \\
\midrule
Synthetic --- Commons Lang & 11.95\% & 5.85\%  & $-$51\% \\
Synthetic --- Commons IO   & 13.49\% & 6.97\%  & $-$48\% \\
Synthetic --- Guava        & 11.76\% & 5.27\%  & $-$55\% \\
\midrule
Real inference --- Commons Lang & 14.70\% & 10.07\% & $-$31\% \\
\bottomrule
\end{tabular}
\end{table}

The synthetic-spec reduction is larger than the real-inference reduction because synthetic specs are dominated by exactly the patterns the optimisation targets: every spec carries the default-omittable \texttt{assignable \textbackslash nothing} (optimisation~2), every \texttt{requires} clause is of the form \texttt{p != null} (optimisation~3 token), and every \texttt{ensures} clause references \texttt{\textbackslash result} (also a token). The real-inference corpus contains a wider variety of clause shapes; many strings still benefit from the dictionary but a long-tail of arithmetic expressions, range comparisons, and quantified subexpressions are not covered by the current twelve tokens. The real-inference reduction is therefore a more realistic estimate of what a production deployment should expect.

The v1 format remains readable; the reader supports both shapes so embedded JARs produced before the format change continue to load.

\subsection{Threats specific to this measurement}

\begin{itemize}
    \item \emph{Three libraries plus self-test}, not the protocol's full ten external libraries. The fidelity analysis above identifies the residual risks (interface methods, multi-release JARs); on libraries that exercise either pattern the embedder reaches 100\%. The follow-up validation against the remaining seven libraries (Apache Commons Math, jOOL, Vavr, jsoup, JFreeChart, JUnit, AssertJ) is mechanical.
    \item \emph{JDK 21 only.} Older JDKs (8, 11, 17) and ahead-of-time compilers (GraalVM \texttt{native-image}) are deferred. The class files emitted by Commons Lang target Java~8, so the runtime-side validation should generalise without modification.
    \item \emph{Source-to-bytecode key resolution is name+arity only.} 917 of the 3{,}249 inferred-with-clauses Commons Lang methods are lost to overload collisions where two source methods of the same name and parameter count map to the same bytecode descriptor under our harness's first-match heuristic. A full type-resolving harness (parsing each source parameter's Java type and converting to its JVM descriptor) would close this gap; the present harness already demonstrates 100\% fidelity on every spec it embeds, so the residual is a measurement limitation rather than an embedder limitation.
\end{itemize}

\section{Discussion}
\label{sec:discussion}

\subsection{Why an annotation, not a class-file attribute?}

Custom class-file attributes (JVMS~\S4.7.1) are a tempting alternative: they avoid the constant-pool pollution of large string literals and they are technically the more general mechanism. Three reasons argue against them in practice:

\begin{itemize}
    \item \emph{Tool visibility.} Standard reflection ignores custom attributes. IDEs do not display them. ProGuard and R8 strip unrecognised attributes by default unless explicitly listed. Each downstream consumer would need a custom reader.
    \item \emph{Repeatable annotations are no longer onerous.} Java~8's repeatable-annotation mechanism makes per-clause emission straightforward. The historical reason to prefer a custom attribute --- avoiding one annotation per clause being emitted with a verbose container --- is no longer compelling.
    \item \emph{Forward compatibility.} An annotation-based carrier with a \texttt{Kind} enum allows new clause kinds to be added without breaking existing readers. A custom attribute's binary layout is harder to evolve compatibly.
\end{itemize}

We therefore recommend the annotation carrier as the primary form. The custom-attribute alternative remains usable for cases where the annotation is unsuitable (signed JARs, aggressive obfuscation), and our sidecar JAR fills the same gap with even less coupling to the bytecode.

\subsection{Generality}

Our format design is independent of the inferrer that produces the specifications. Any source of JML clauses --- the inferrer of~\cite{edmonds2026inference}, Daikon~\cite{ernst2007daikon}, hand-written contracts, contracts produced by a future LLM-assisted refinement step --- can be embedded by the same mechanism. The format also generalises to languages other than Java that target the JVM (Kotlin, Scala) at the cost of those languages providing their own translators between source-level annotations and the bytecode \texttt{@JmlSpecs} representation.

For non-JVM languages, the OpenAPI extension carries equivalent information at the service boundary. The bytecode side is JVM-specific; the service-boundary side is language-agnostic.

\subsection{Limitations}

The embedder is consumer-tolerant but the consumer has work to do. A reader that wants to act on the embedded clauses must (a) parse the canonical-form JML strings, (b) substitute parameter names, and (c) decide what to do with the result. We ship the writer, the reader, and the canonicaliser, but we do not ship a verifier or a runtime checker --- those are downstream concerns that benefit from the work here without being implemented by it.

We note one substantive limitation: roundtrip equality is on the canonical-form string, not on a normalised semantic theory, so equivalent specifications written in different forms will appear distinct on the wire. We chose this separation of concerns deliberately in Section~\ref{sec:format}; semantic normalisation is a separable layer.

\section{Threats to Validity}
\label{sec:threats}

\paragraph{Internal validity.}
The 100\% roundtrip-fidelity figure depends on the writer covering every method shape that ASM emits. Section~\ref{sec:empirical} discusses the original 3\% gap that arose when the writer's lazy emission missed the abstract / interface case (only \texttt{visitEnd} fires for them). Adding \texttt{visitEnd} to the trigger set closed the gap on the libraries tested. The same risk pattern applies to method shapes specific to future JVM features --- e.g., abstract record components, hidden classes (JEP~371) --- which would need their own triggers if they exhibit different visit-call orderings. The mitigation is the same: a corpus-scale validation per JVM release.

The byte-overhead figure is measured under a deliberately bulky synthetic-spec workload (one annotation per parameter for every method). Real inferred specifications are typically smaller per method, so the reported 11.76\%--13.49\% range should be read as an upper bound for the embedder's contribution to JAR size. The 13.49\% figure on Commons IO reflects its higher proportion of interfaces (each abstract method receives an embedded annotation just like a concrete method); a workload-aware embedder that elides annotations on abstract methods whose specs are pure pass-throughs would lower the figure further.

\paragraph{External validity.}
The two libraries reported are mature, well-structured, and ship through Maven Central. Generalisation to less-tidy code (older artefacts, vendored dependencies, hand-tuned ProGuard-stripped JARs) will need separate validation. The negative test confirms that the consumer-side contract (load classes without the annotation type) holds at JDK 21; older JDKs (8, 11, 17) and ahead-of-time compilation paths (GraalVM \texttt{native-image}) are deferred.

The OpenAPI extension is validated by construction (it is a JSON Schema valid against OpenAPI 3.x's extension rules). Empirical validation against a real microservice is part of the planned follow-up evaluation but does not appear in the present article.

\paragraph{Construct validity.}
"Roundtrip fidelity" is operationalised as set-equality on \texttt{(internalName, methodName, descriptor)} keys. Two methods with identical signatures but different bytecode bodies (for example, the same constructor compiled with different optimisation levels) appear identical to the harness. This is the right construct for the embedder's job --- the embedder carries metadata, not bytecode --- but a reader that needs to attach annotations to a specific compilation of a method would need additional disambiguation.

\paragraph{Conclusion validity.}
The empirical claims are point measurements on two libraries; no statistical significance test is meaningful. The claims are robust in the sense that the gap-causing patterns are identified and fixable, and the absolute throughput numbers are well above the thresholds CI/CD workflows require, but a wider corpus (the protocol's full ten libraries) is required before the claims generalise.

\section{Related Work}
\label{sec:related}

\paragraph{Specification languages.}
JML~\cite{jml,leavens2006design} provides the contract vocabulary we transport. Eiffel's design-by-contract~\cite{meyer1992design} originated the broader approach; Spec\#~\cite{barnett2005weakest} ported it to .NET. None of these languages address the distribution problem we tackle --- they assume the consumer has source.

\paragraph{Specification inference.}
The inference engine of~\cite{edmonds2026inference} is the immediate predecessor of this work. Daikon~\cite{ernst2007daikon} infers likely invariants from execution traces; Houdini and its descendants infer candidate annotations and discharge them through a verifier. PreInfer~\cite{preinfer2017} uses dynamic symbolic execution. Jdoctor~\cite{jdoctor2018} extracts exception preconditions from Javadoc. None of these systems address what happens to the inferred specifications after they are produced; we provide the missing distribution layer.

\paragraph{Annotation-based contracts.}
JSR-305 captured nullability annotations and was withdrawn before final standardisation; the surviving informal usage covers only nullability. The Checker Framework's annotation vocabulary is type-system-shaped and does not generalise to the relational and quantified clauses JML supports. Cofoja (Contracts for Java) emits runtime checks from \texttt{@Requires}/\texttt{@Ensures} annotations but does not survive without its annotation processor on the classpath; our work is consumer-tolerant by design.

\paragraph{Bytecode metadata transport.}
ASM~\cite{asm} is the standard library for class-file rewriting and underpins our embedder. ByteBuddy and Javassist provide higher-level APIs over similar primitives. None of these libraries prescribe a metadata format for specifications --- they are mechanism, not policy. JEP~181 (nest-mates) and JEP~371 (hidden classes) extend what the JVM accepts as class-file content but do not change the annotation channel.

\paragraph{API description.}
OpenAPI~3.x is the standard for HTTP service descriptions; gRPC's \texttt{.proto} files cover RPC. Pact~\cite{pact} captures consumer-driven contracts as example interactions; it is example-based rather than logical, so it is a transport for examples rather than for specifications, and we specifically reject example-based encoding (Section~\ref{sec:format}). The OpenAPI extension we propose is the smallest viable addition to an existing widely deployed format.

\paragraph{Generating OpenAPI from source.}
Recent work on generating OpenAPI descriptions from Java bytecode and source code includes Respector, Prophet, and AutoOAS~\cite{Lercher2024-AutoOAS}, which extract REST endpoint metadata from Spring Boot applications. These approaches generate the OpenAPI description itself; our \texttt{x-jml-*} extension is orthogonal to that work --- it augments an OpenAPI description with logical contracts, regardless of how the description was produced.

\paragraph{Mutation testing and test generation.}
The mutation-testing literature is cited in detail in~\cite{edmonds2026inference}. The relationship is: that work measured what specifications do for LLM-generated tests, and the present work makes the specifications transportable to LLMs and other consumers that do not have the source.

\section{Conclusion}
\label{sec:conclusion}

We asked whether automatically inferred specifications can be embedded in compiled Java artefacts and REST API descriptions with sufficient fidelity that downstream consumers can use them as if they had the source. We answered in the affirmative, with two specific constructions:

\begin{enumerate}
    \item A runtime-retained Java annotation, \texttt{@JmlSpecs}, written into class-file \texttt{RuntimeVisibleAnnotations} attributes by an ASM-based embedder. Roundtrip fidelity is 100.0\% on three real-world libraries (Apache Commons Lang, Apache Commons IO, Guava --- totalling 20{,}546 methods) and on a real-inferred-specs corpus of 2{,}332 Commons Lang specs spanning quantifiers, \texttt{\textbackslash old} references, and branch-conditional implications. Bytecode-size overhead after the format-optimisation step described in Section~\ref{subsec:format-opt} is 5.27--6.97\% on synthetic specs and 10.07\% on real inferred specs (a 48--55\% / 31\% reduction over the initial per-clause-annotation format). Embedding throughput is in the 12{,}000--22{,}000 methods/second range; reading throughput is in the 119{,}000--328{,}000 methods/second range.

    \item An OpenAPI 3.x extension family (\texttt{x-jml-requires}, \texttt{x-jml-ensures}, \texttt{x-jml-assignable}, \texttt{x-jml-signals}, \texttt{x-jml-invariant}) per-operation and per-schema, plus a sidecar JSON fallback for documents generated by frameworks. The extension is preserved by standard OpenAPI parsers without modification.
\end{enumerate}

Our negative test confirms that classes embedded with \texttt{@JmlSpec} load correctly in JVMs that have no awareness of the annotation type --- the embedder is consumer-tolerant by design. Together with the inference engine of~\cite{edmonds2026inference}, the path from source code to consumer-accessible contract is complete: the inferrer produces specifications, the embedder ships them alongside compiled artefacts, and the OpenAPI extension carries them across service boundaries.

We see three threads of follow-up work. First, validation across the remaining seven libraries identified in our protocol (Apache Commons Math, jOOL, Vavr, jsoup, JFreeChart, JUnit, AssertJ). Second, integration with the LLM-based test generator of~\cite{edmonds2026inference}, which can now consume specifications directly from the JAR rather than requiring source. Third, an extension of the inference itself to compose specifications across method boundaries; the carrier built here is what makes that composition transportable to consumers.

The replication package, including the full source of the embedder/extractor pair, the OpenAPI extension JSON Schema, and the empirical validation harness, accompanies this paper.


\newpage
\bibliographystyle{IEEEtran}
\bibliography{../article1/references}

\end{document}
